\documentclass[french]{article}

%Chargement des paquets
\usepackage{amsmath}
\usepackage{amsfonts}
\usepackage{amssymb}
\usepackage{enumerate}
\usepackage{mathtools}
\usepackage{bbm}
\usepackage{xparse, etoolbox}
\usepackage{enumerate}
\usepackage{enumitem}
\usepackage{mathabx}
\usepackage{babel}[french]

%environment
\usepackage{amsthm}
\usepackage{keytheorems}
\usepackage{hyperref} 

%Interface théorème
\newkeytheorem{lem}[
	name = Lemme
]

\newkeytheorem{prop}[
	name = Proposition
]

\newkeytheorem{thm}[
	name = Théorème
]

\newkeytheorem{coro}[
	name = Corollaire
]

\renewcommand*{\proofname}{Démonstration}

\theoremstyle{definition}
\newtheorem{defn}{Définition}

\theoremstyle{definition}
\newtheorem*{nota}{Notation}

\theoremstyle{definition}
\newtheorem*{limi}{Liminaire}

\theoremstyle{remark}
\newtheorem{rmq}{Remarque}

\theoremstyle{plain}
\newtheorem*{fait}{Fait}


%Ensembles de nombres
\newcommand{\N} {\mathbb{N}}
\newcommand{\Ne}{\N^\ast}
\newcommand{\Z} {\mathbb{Z}}
\newcommand{\Q} {\mathbb{Q}}
\newcommand{\R} {\mathbb{R}}
\newcommand{\Rb}{\overline{\mathbb{R}}}
\newcommand{\Rp}{\R_+}
\newcommand{\Rm}{\R_-}
\newcommand{\K} {\mathbb{K}}
\newcommand{\Ket} {\mathbb{K^{\ast}}}
\newcommand{\Cx}{\mathbb{C}}

%Ensembles, tribus
\newcommand{\A}{\mathbb{A}}
\renewcommand{\P}{\mathcal{P}}
\newcommand{\T}{\mathcal{T}}
\newcommand{\F}{\mathcal{F}}
\newcommand{\C} {\mathcal{C}}
\newcommand{\ouv}{\mathcal{O}}
\newcommand{\B}{\mathcal{B}}
\newcommand{\M}{\mathcal{M}}
\newcommand{\etag}{\mathcal{E}}
\newcommand{\etagOp}{\etag^{+}(\Omega)}
\newcommand{\etagO}{\etag(\Omega)}
%%Lp avant quotientage
\newcommand{\Lic}{\mathcal{L}^1}
\newcommand{\Lpc}{\mathcal{L}^p}
\newcommand{\Linfc}{\mathcal{L}^{\infty}}
\newcommand{\Lifull}{\Lic(\Omega, \B, m)}
\newcommand{\Lpfull}{\Lpc(\Omega, \B, m)}
\newcommand{\Linffull}{\Linfc(\Omega, \B, m)}
%%Lp après quotientage
\newcommand{\Li}{L^1}
\newcommand{\Ld}{L^2}
\newcommand{\Lp}{L^p}
\newcommand{\Linf}{L^{\infty}}
\newcommand{\Listd}{\Li(\Omega, m)} 
\newcommand{\Ldstd}{\Ld(\Omega, m)} 
\newcommand{\Lpstd}{\Lp(\Omega, m)} 

%Quelques opérateurs
\DeclareMathOperator{\codim}{codim}
\DeclareMathOperator{\rg}{rg}
\DeclareMathOperator{\tr}{tr}
\DeclareMathOperator{\vect}{Vect}
\DeclareMathOperator{\ima}{Im}
\DeclareMathOperator{\id}{Id}
\DeclareMathOperator{\sign}{sign}
\DeclareMathOperator{\ext}{ext}
\newcommand{\ind}{\mathbbm{1}}
\newcommand*{\spr}[2]{\left(\,#1\,\middle|\,#2\,\right)}
\newcommand{\interior}[1]{%
  {\kern0pt#1}^{\mathrm{o}}%
}
\DeclareMathOperator{\card}{card}
\DeclareMathOperator{\ppcm}{ppcm}

%Commandes de pure flemme
\newcommand{\veps}{\varepsilon}
\newcommand{\intab}{\int_{a}^{b}}
\newcommand{\intR}{\int_{-\infty}^{\infty}}
\newcommand{\intinf}{\int_{- \infty}^{+ \infty}}
\newcommand{\equi}{\Leftrightarrow}
\newcommand{\Kev}{$\K$-espace vectoriel }
\newcommand{\Kevs}{$\K$-espaces vectoriels }
\newcommand{\Rev}{$\R$-espace vectoriel }
\newcommand{\Revs}{$\R$-espaces vectoriels }
\newcommand{\Cev}{$\Cx$-espace vectoriel }
\newcommand{\Cevs}{$\Cx$-espaces vectoriels }
\newcommand{\evn}{(E, \norm{.})}
\newcommand{\interf}[2]{[#1,#2]}
\newcommand{\limb}{\lim\limits_}
\newcommand{\lebn}{\lambda_n}
\newcommand{\pp}{\text{~presque partout}}
\newcommand{\derivp}[2]{\frac{\partial #1}{\partial #2}}
\newcommand{\famiu}{(u_i)_{i \in I}}
\newcommand{\sev}{\text{~s.e.v.}}
\newcommand{\Mnp}{M_{n,p}(\K)}
\newcommand{\Mn}{M_{n}(\K)}
\newcommand{\tq}{\text{~t.q.~}}
\newcommand{\mq}{~montrer~que~}
\newcommand{\et}{\quad \text{et} \quad}
\newcommand{\ou}{\text{~ou~}}
\newcommand{\Kx}{\K[X]}


%Commandes de gars chelou
\renewcommand{\subset}{\subseteq}

%Commande restriction, ty duckduckgo
\def\restr#1#2{\mathchoice
              {\setbox1\hbox{${\displaystyle #1}_{\scriptstyle #2}$}
              \restraux{#1}{#2}}
              {\setbox1\hbox{${\textstyle #1}_{\scriptstyle #2}$}
              \restraux{#1}{#2}}
              {\setbox1\hbox{${\scriptstyle #1}_{\scriptscriptstyle #2}$}
              \restraux{#1}{#2}}
              {\setbox1\hbox{${\scriptscriptstyle #1}_{\scriptscriptstyle #2}$}
              \restraux{#1}{#2}}}
\def\restraux#1#2{{#1\,\smash{\vrule height .8\ht1 depth .85\dp1}}_{\,#2}} 

%Quelques normes
\DeclarePairedDelimiter\abs{\lvert}{\rvert}
\DeclarePairedDelimiter\labs{\bigl\lvert}{\bigr\rvert}
\DeclarePairedDelimiter\norm{\lVert}{\rVert}
\DeclarePairedDelimiterXPP\normi[1]{}\lVert\rVert{_1}{\ifblank{#1}{\:\cdot\:}{#1}}
\DeclarePairedDelimiterXPP\normd[1]{}\lVert\rVert{_2}{\ifblank{#1}{\:\cdot\:}{#1}}
\DeclarePairedDelimiterXPP\normp[1]{}\lVert\rVert{_p}{\ifblank{#1}{\:\cdot\:}{#1}}
\DeclarePairedDelimiterXPP\normq[1]{}\lVert\rVert{_q}{\ifblank{#1}{\:\cdot\:}{#1}}
\DeclarePairedDelimiterXPP\norminf[1]{}\lVert\rVert{_\infty}{\ifblank{#1}{\:\cdot\:}{#1}}

%Bracket en angle
\DeclarePairedDelimiter\angl{\langle}{\rangle}

%Set, parenthèses
\newcommand{\set}[1]{\{ #1 \}}
\newcommand{\bigset}[1]{\bigl\{ #1 \bigr\}}
\newcommand{\Bigset}[1]{\Bigl\{ #1 \Bigr\}}
\newcommand{\bigpar}[1]{\bigl( #1 \bigr)}
\newcommand{\Bigpar}[1]{\Bigl( #1 \Bigr)}

%Opitmisation des opérateurs de comparaison
\DeclarePairedDelimiter\tio{o (}{)}
\DeclarePairedDelimiter\gro{O (}{)}


\begin{document}

Dans ce document, $I$ désigne un intervalle de $\R$ (non réduit à un point) et $\K$ désigne $\R$ ou $\Cx$.
On prendra, le plus souvent, $\N$ tout entier comme support des suites (et des séries) mais on peut étendre les résultats à des supports 
$S \subset \N$ infinis.

\section{Convergence simple et uniforme}

\begin{defn}
Soit $(f_n)_\N$ une suite de fonctions de $I$ dans $\K$.
Soit $(S_N)_\N$, la suite de fonctions définie par :
\[ \forall N \in \N, \forall x \in I, S_N(x) = \sum_{n=0}^N f_n(x) . \]
On appelle $(S_N)_\N$, la série de fonctions de terme général $f_n$, et on la note $\sum f_n$.
\end{defn}

\begin{defn}[Convergence simple]
Soit $(f_n)_\N$ une suite de fonctions de $I$ dans $\K$.
Lorsque la suite $(S_N)_\N$ des sommes partielles converge simplement vers une fonction $S$, on dit que la série $\sum f_n$ converge 
simplement et on note $S = \sum_{n=0}^{\infty}f_n$.
\end{defn}

\begin{defn}[Convergence absolue]
Soit $(f_n)_\N$ une suite de fonctions de $I$ dans $\K$.
On dit que la série $\sum f_n$ converge absolument si la série $\sum \abs{f_n}$ converge simplement.
\end{defn}

\begin{thm}
Pour toute série de fonctions de $I$ dans $\K$, la convergence absolue implique la convergence simple.
\end{thm}

\begin{proof}
Pour $x \in I$ fixé, la convergence absolue de la série numérique $\sum f_n(x)$ implique la convergence de cette série.
Donc, la convergence absolue de la série de fonctions implique la convergence des séries numériques $\sum f_n(x)$ pour toutes les
valeurs de $x \in I$, ce qui est la définition de la convergence simple.
\end{proof}

\begin{rmq}
Si une série de fonctions converge simplement mais pas absolument, elle est dite \og semi-convergente \fg.
\end{rmq}

\begin{defn}[Convergence uniforme]
Soit $(f_n)_\N$ une suite de fonctions de $I$ dans $\K$.
Lorsque la suite $(S_N)_\N$ des sommes partielles converge uniformément vers une fonction $S$, 
on dit que la série $\sum f_n$ converge uniformément.
\end{defn}

\begin{thm}
Pour toute série de fonctions de $I$ dans $\K$, la convergence uniforme implique la convergence simple.
\end{thm}

\begin{proof}
C'est un résultat classique des suites de fonctions.
\end{proof}

\begin{rmq}
Si on note
\[ R_N = \sum_{n=N+1}^{\infty} f_n , \]
il y a alors équivalence entre convergence uniforme de la série $\sum f_n$ et convergence uniforme de la suite des restes $(R_N)_\N$ 
vers la fonction nulle.
\end{rmq}

\section{Convergence normale}

\begin{defn}[Convergence normale]
Soit $(f_n)_\N$ une suite de fonctions de $I$ dans $\K$.
On dit que la série $\sum f_n$ converge normalement si la série $\sum \norminf{f_n}$ est convergente.
\end{defn}

\begin{rmq}
En pratique, une série de fonction $\sum f_n$ converge normalement si, et seulement si, il existe une suite réelle $(u_n)_\N$ telle que :
\begin{enumerate}[label=(\roman*)]
\item $\forall n \in \N, \forall x \in I, \abs{f_n(x)} \leq u_n$ :
\item $\sum u_n$ converge.
\end{enumerate}
\end{rmq}

\begin{thm}
Soit $(f_n)_\N$ une suite de fonctions de $I$ dans $\K$.
Si la série $\sum f_n$ converge normalement alors elle converge absolument et uniformément.
\end{thm}

\begin{proof}
Il est clair que $\forall x \in I, \abs{f_n(x)} \leq \norminf{f_n}$ donc si $\sum \norminf{f_n}$ converge alors, la série $\sum f_n$ converge
absolument. \\
On note : 
\[ \forall x \in I, S(x) = \sum_{n=0}^{\infty} f_n(x) \et S_N(x) = \sum_{n=0}^N f_n(x) ,\]
On vient de démontrer que si la série $\sum f_n$ converge normalement, alors elle converge absolument (et donc simplement), ce qui
implique que $S$ est bien définie (comme limite simple de la suite $(S_N)_\N$). Or :
\begin{align*}
\forall x \in I, \abs{S(x) - S_N(x)} & = \abs*{ \sum_{n=N+1}^{\infty} f_n(x) } \\
& \leq  \sum_{n=N+1}^{\infty}  \abs{f_n(x)} \\
& \leq \sum_{n=N+1}^{\infty} \norminf{f_n}
\end{align*}
Donc $\norminf{S(x) - S_N(x)} \leq \sum_{n=N+1}^{\infty} \norminf{f_n} \to 0$ pour $N \to \infty$.
\end{proof}

\begin{thm}

\end{thm}

\begin{proof}

\end{proof}


\end{document}